% Agent 25: Pages 439--442 (PDF pp.48--51)
% Sections 5.9 (cont.) inner region profiles
%          5.10 Spectrum of Radiation from Disk
%          5.11 Heating of Outer Region by X-Rays from the Inner Region
% =============================================================================

%% ----------------------------------------------------------------
%% 5.9 (continued) — Radial profiles
%% ----------------------------------------------------------------

and the characteristic timescale for the gas to move inward from radius $r$ to the
inner edge of the disk
%
\begin{equation}
\tag{5.9.5}
\Delta(r) = -r/v^r\,.
\end{equation}

\noindent
$\Delta(r) = (2~\text{sec})(\alpha^{-4/5} M_*^{3/10})_*^{5/4} \mathscr{g}^{1/10} \mathscr{g}^{-4/5} \mathscr{g}^{1/2} \mathscr{g}^{-7/20}$

\noindent
$\times\,\mathscr{g}^{-1/20}\,\mathscr{g}^{7/10}$.

It is worth noting that the relativistic correction $\mathscr{Q}$ goes to zero smoothly at the
inner edge of the disk
%
\[
\mathscr{Q} \to 0 \quad \text{as } r \to r_{ms}\,.
\]

The transition to the middle region occurs where $\tau_{ff}/\tau_{es} \sim 1$---i.e., at
%
\begin{equation}
\tag{5.9.7}
r_* = r_{om*} \equiv 2 \times 10^3 (M_*^{-2/3} \dot{M}_{0*}^{2/3})_* \mathscr{g}^{2/3} \mathscr{g}^{-8/15} \mathscr{g}^{-1/3}\,\mathscr{g}^{-1/3}\,\mathscr{Q}^{2/3}\,.
\end{equation}
%
$\simeq 100$ for ``supermassive case.''

\medskip
\noindent
\textit{Middle region}

\noindent
$p = p^{(\text{gas})}$, $\bar{\kappa} = \kappa_{es}$.  In this region the equations of structure (5.8.1) and (5.6.14)
yield:
%
\begin{align}
F &= (0.6 \times 10^{26}~\text{erg/cm}^2~\text{sec})(M_*^{-2} \dot{M}_{0*})_*^{-3}\,\mathscr{g}^{-1}\,\mathscr{g}^{-1/2}\,\mathscr{Q}\,,
\notag\\
\Sigma &= (9 \times 10^2~\text{cm})(M_*^{3/4}\dot{M}_{0*}^{-3/20})_*^{-8}\,\mathscr{g}^{-9/2}\,\mathscr{g}^{-11/10}\mathscr{g}^{1/2}
\notag\\
&\quad \times\,\mathscr{g}^{-33/40}\,\mathscr{g}^{-19/40}\,\mathscr{Q}^{-23/20}\,,
\notag\\
h &= (9 \times 10^2~\text{cm})(M_*^{3/4}\dot{M}_{0*}^{-3/20})_*^{-8}\,\mathscr{g}^{-9/2}\,
\mathscr{g}^{-11/10}\mathscr{g}^{1/2}
\notag\\
&\quad \times\,\mathscr{g}^{-33/40}\,\mathscr{g}^{-19/40}\,\mathscr{Q}^{-23/20}\,,
\notag\\
\rho_0 &= (8 \times 10^1~\text{g/cm}^3)(\alpha^{-7/10} M_*^{5/4} \dot{M}_{0*}^{1/20})_*^{-15/8}\,\mathscr{g}^{-17/20}\mathscr{g}^{3/10}
\notag\\
T &= (8 \times 10^7~\text{K})(\alpha^{-1/5}M_*^{-1/2}\dot{M}_{0*}^{1/0})_*^{-3/4}\,\mathscr{g}^{-1/10}\mathscr{g}^{-1/5}\mathscr{g}^{-3/20}
\notag\\
&\quad \times\,\mathscr{g}^{1/20}\mathscr{g}^{3/10}\,,
\notag\\
\tau_{ff} &= (2 \times 10^3)(\alpha^{-4/5}M_*^{5/4}\mathscr{g}^{-2/5}\mathscr{g}^{1/5}\mathscr{g}^{1/2}\,\mathscr{g}^{-3/5}\,\mathscr{g}^{1/5}\,\mathscr{Q}^{1/5}\,,
\tag{5.9.6}\\
B &\leq (7 \times 10^8~\text{G})(\alpha^{1/20}M_*^{-7/8}\dot{M}_{0*}^{17/40})_*^{-21/16}\,\mathscr{g}^{-19/40}\mathscr{g}^{1/20}\mathscr{g}^{-17/80}
\notag\\[6pt]
\left(\frac{p^{(\text{gas})}}{p^{(\text{rad})}}\right) &= 4(\alpha^{-1/10} M_*^{1/4} \dot{M}_{0*}^{-7/20})_*^{3/2}\,\mathscr{g}^{-3/5}\,\mathscr{g}^{1/5}\,\mathscr{g}^{1/2}\,\mathscr{g}^{-4/5}\,,
\notag\\[4pt]
\left(\frac{\tau_{ff}}{\tau_{es}}\right) &= (0.02)(\alpha^{-1/10} M_*^{1/10} M_*^{1/5})_*^{21/10}\,\mathscr{g}^{9/5}\mathscr{g}^{2/5}\,\mathscr{g}^{1/2}\,\mathscr{Q}^{-4/5}\,,
\notag\\[4pt]
r_*^{} &= (0.6 \times 10^{-5})(M_* \dot{M}_{0*})*^{3/2}\,\mathscr{g}^{\,}\mathscr{g}^{2/1}\,\mathscr{g}^{1/2}\,\mathscr{Q}^{-1}\,,
\notag\\
&= 50(\alpha^{-4/5} M_*^{-9/10} \dot{M}_{0*}^{1/10})_*^{3/2}\,\mathscr{g}^{-4/5}\,\mathscr{g}^{1/2}\,\mathscr{g}^{1/2}\mathscr{g}^{-3/10}\,\mathscr{Q}^{3/5}\,.
\notag
\end{align}
%
$\Delta(r) = (0.7~\text{sec})(\alpha^{-4/5} M_*^{5/8} \dot{M}_{0*}^{1/5})_*^{21/5}\,\mathscr{g}^{-4/5}\,\mathscr{g}^{1/2}\,\mathscr{g}^{-3/10}\,\mathscr{Q}^{3/5}$.

\begin{equation}
\tag{5.9.8}
\end{equation}

The transition from the middle region to the inner region occurs where
$p^{(\text{gas})}/p^{(\text{rad})} \sim 1$---i.e., at
%
\begin{equation}
\tag{5.9.9}
r_* = r_{mi*} \equiv 40(\alpha^{2/21} M_*^{-2/3} \dot{M}_{0*}^{16/21})
\mathscr{g}^{20/21}\mathscr{g}^{-36/21}\,\mathscr{g}^{-8/21}\,\mathscr{g}^{-10/21}\,\mathscr{Q}^{16/21}
\end{equation}
%
$\simeq 4$ for supermassive case.

Thus, in the supermassive case the ``middle region'' extends all the way---or
almost all the way---into the inner edge of the disk.

\medskip
\noindent
\textit{Inner region}

\noindent
$p = p^{(\text{rad})}$, $\bar{\kappa} = \kappa_{es}$.  In this region the equations of structure (5.8.1) and (5.6.14)
yield:
%
\begin{align}
F &= (0.6 \times 10^{26}~\text{erg/cm}^2~\text{sec})(M_*^{-2} \dot{M}_{0*})_*^{-3}\,\mathscr{g}^{-1}\,\mathscr{g}^{-1/2}\,\mathscr{Q}\,,
\notag\\
\Sigma &= (20~\text{g/cm}^2)(\alpha^{-1}M_*\dot{M}_{0*})_*^{3/2}\,\mathscr{g}^{-3/2}\,\mathscr{g}^{3}\,\mathscr{g}^{1/2}\,\mathscr{g}^{-1/2}\,\mathscr{Q}^{-1}\,,
\notag\\
h &= (1 \times 10^5~\text{cm})(\dot{M}_{0*}\mathscr{g}^{-3}\,\mathscr{g}^{-1/2}\,\mathscr{g}^{-1}\,\mathscr{Q}\,,
\notag\\
\rho_0 &= (1 \times 10^{-4}~\text{g/cm}^3)(\alpha^{-1}M_*\dot{M}_{0*})*^{3/2}\,\mathscr{g}^{-4}\mathscr{g}^{2}\,\mathscr{Q}^{-2}\,,
\tag{5.9.10}\\
T &= (4 \times 10^7~\text{K})(\alpha^{-1/4}M_*^{-1/4})_*^{-9/8}\,\mathscr{g}^{-1/2}\mathscr{g}^{1/2}\,\mathscr{Q}^{1/4}\,,
\notag\\
\tau_{es} &= 8(\alpha^{-1}M_*\dot{M}_{0*}^{1/2})_*^{3/2}\,\mathscr{g}^{-3/2}\,\mathscr{g}^{1/2}\,\mathscr{g}^{-1/2}\,\mathscr{g}^{-1}\,\mathscr{Q}^{-1}\,,
\notag\\
B &\leq (7 \times 10^7~\text{G})(M_*\dot{M}_{0*}^{1/2})_*^{-3/4}\,\mathscr{g}^{-1/2}\,\mathscr{g}^{1/2}\,\mathscr{Q}^{-1}\,,
\notag
\end{align}
%
\[
\left(\frac{p^{(\text{gas})}}{p^{(\text{rad})}}\right)
= (5 \times 10^{-5})(\alpha^{-1/4}M_*^{1/4}M_{0*}^{2})_*^{1/8}\,\mathscr{g}^{-5/2}\mathscr{g}^{6/4}\,\mathscr{Q}^{-2}\,,
\]
%
\[
r^{**} \equiv (\tau_{es}\tau_{ff})^{1/2}
= (5 \times 10^{-5})(\alpha^{-17/16}\,M_*^{1/16}\,\mathscr{g}^{2/5}\,\mathscr{g}^{1/2}\,\mathscr{Q}^{-2/5}\,\mathscr{g}^{-25/16}
\]
%
\[
= (2 \times 10^{-3})(\alpha^{-17/16}\,M_*^{31/16}\,\dot{M}_{0*}^{3/1})
\,\mathscr{g}^{93/32}\,\mathscr{g}^{-25/8}\mathscr{g}^{1/2}\,\mathscr{g}^{25/16}
\]

\[
\Delta(r) = (2 \times 10^{-4}~\text{sec})(\alpha^{-1}M_*^{3/8}\dot{M}_{0*}^{3})_*^{2}\,\mathscr{g}^{-25/8}\,\mathscr{g}^{1/2}\,\mathscr{Q}^{1/2}\,\mathscr{g}^{-2}\,.
\]

Notice that for $\mathscr{Q} \leq 1$ ``typical'' cases of interest (galactic X-ray sources with
$M_* \simeq M_{0*} \sim 1$, $a \simeq M$ (``maximal
Kerr''), the disk will be optically thin over a narrow region between $r \sim 2M$ and
$r \sim 10M$.  Inside this region it will be thick (because $\mathscr{Q} \to 0$ as $r \to M$); outside it
will also be thick.

In any case where the model
is everywhere thin in the sense that $h/r \leq 0.1$.  In any case where $M_* \sim 10^7$, $\dot{M}_{0*} \sim 10^5$ the disk
predicts $h/r \gtrsim 1$, the model is self-inconsistent.

%% ================================================================
%% 5.10  Spectrum of Radiation from Disk
%% ================================================================

\subsection{Spectrum of Radiation from Disk}
\label{sec:5.10}

\noindent
\textit{Outer region}

In the outer region electron-scattering opacity is negligible compared to free-free
opacity, and the disk is optically thick.  Consequently, the disk's surface
is able to build up a blackbody spectrum (see \S\ref{sec:2.6}).  The ``surface temperature''
of the disk, which characterizes the blackbody spectrum, is
%
\begin{equation}
\tag{5.10.1}
T_s = \left(\frac{4F}{b}\right)^{1/4}
= (3 \times 10^5~\text{K})(M_*^{-1/2}\dot{M}_{0*}^{1/4})_*^{-3/4}\,
\mathscr{g}^{-1/4}\,\mathscr{g}^{-1/8}\,\mathscr{Q}^{1/4}
\end{equation}
%
\[
\simeq (1 \times 10^5~\text{K})(M_*^{4}(r_*/r_{om*})^{-3/4}\,.
\]

Here $r_{om*}$ is the inner edge of the outer region [eq.\ (5.9.7)].  Thus, the outer
region of the disk's surface emits a blackbody spectrum with temperature
${\sim}\,10^5$\,K to ${\sim}\,10^5$\,K.  For typical galactic X-ray sources ($M_* \sim M_{0*} \sim 1$) the
inner edge of the outer region is located at $r_{om} \simeq 10^9$\,cm---compared to a size
$r \sim 3 \times 10^{11}$\,cm for the inner Roche lobe of the disk region, and compared to a size
$r \sim 10^6$\,cm for the inner edge of the disk if the central object is a black
hole or neutron star, and $r \sim 10^9$\,cm if it is a white dwarf.

\medskip
\noindent
\textit{Middle and inner region}

In the middle region and inner region electron-scattering opacity modifies the
emitted spectrum so it is no longer blackbody.  An unsophisticated way to calculate
the modification is this: (i)~At some fixed radius consider all photons that
emerge from the disk with some fixed frequency~$\nu$.  After being formed by free-free
emission or some process, these photons must ``random-walk'' their
way through scattering electrons before they can emerge from the disk's
surface.  Let $y_\nu$ be the depth in the disk at which these photons are formed, as
measured in g/cm$^2$.
%
\[
y_\nu = \int_{\text{formation point}}^{\text{surface}} \rho_0 \, dz\,.
\]

The mean-free path of the photons between scatterings, as measured in these
same units, is $\lambda = 1/\kappa_{es} = 2.5~\text{g/cm}^2$ (independent of photon frequency).  The
total depth travelled, $y_\nu$, is the product of this mean-free path with the square-root
of the number of scatterings $N_{ss}$
%
\[
y_\nu = \lambda N_{ss}^{1/2} = \kappa_{es}^{-1} N_{ss}^{1/2}
\]
%
(``standard square-root factor for random walks'').  Hence, the total number of
scatterings between emission and emergence from disk is
%
\[
N_{ss} = (\kappa_{es} y_\nu)^2\,.
\]

The depth of the emission point $y_\nu$ is determined by the demand that the total
free-free optical depth traversed along a photon's tortured trajectory be unity:
%
\[
1 = \kappa_{ff}^{\nu}(y_\nu \lambda) = \kappa_{ff}^{\nu}(y_\nu / \kappa_{es}^{1/2})\,.
\]

Here $\tau_{ff}^{\nu}$ like $\tau_{es}$ is measured not along the tortured random-walk photon path,
but rather along the straight-line path of standard radiative transfer theory (\S\ref{sec:2.6}).
Let us summarize: in an atmosphere dominated by electron scattering, those
photons of frequency $\nu$ that escape from the surface are generated at a depth
[Zel'dovich and Shakura (1969), Shakura (1972)]
%
\begin{equation}
\tag{5.10.2}
\tau_{es*}(\text{emission pt.})
= [\tau_{ff}^{\nu}(\text{emission pt.})\,\tau_{es}(\text{emission pt.})]^{1/2}
= (\kappa_{ff}^{\nu}\kappa_{es})^{1/2}\,y_\nu = 1\,.
\end{equation}

[Note: throughout this discussion we have ignored the tiny changes in photon
frequency at each scattering; see \S\ref{sec:2.6} and see below.]

Because $\kappa_{ff}^{\nu} \propto (1 - e^{-x})/x^3$ for $x \gg 1$, and $\simeq x^{-3}$ for $x \ll 1$,
where $x \equiv h\nu/T$,
%
\begin{equation}
\tag{5.10.3}
\kappa_{ff}^{\nu} \; \text{decreases with increasing photon frequency.}
\end{equation}

[eq.\ (2.6.25)], the high frequency part of the spectrum gets formed at greater
depths than the low-frequency spectrum.

At the formation point, the specific intensity will be blackbody, $I_\nu = B_\nu$.  The
specific flux that crosses outward through the surface at depth $\tau_\nu$ (ignoring for
the moment the flux that crosses back inward) thus has the blackbody form
%
\[
F_\nu(\tau_\nu) = \int_0^{\pi/2} B_\nu \cos\theta \, d\Omega = 2\pi B_\nu\,.
\]

But of all the photons in this specific flux, only a fraction $1/(N_{ss})^{1/2}$ ever
reaches the surface of the disk (``standard $\sqrt{N}$ random-walk factor'').  All the
rest eventually get scattered back to depths greater than $y_\nu$, and eventually get
absorbed there.  Hence, the specific flux emerging from the surface of the disk
will be
%
\begin{equation}
\tag{5.10.4}
F_\nu = 2\pi B_\nu (\tau_{es}(\text{at pt.\ where } \tau_{*} = 1))^{-1}\,.
\end{equation}

This is the ``modified spectrum'' which emerges from the middle and inner parts
of the disk.

However, a more sophisticated derivation, with an atmosphere in which tempera\-ture
and density (and hence $\kappa_{ff}^{\nu}$) vary with height, yields essentially the same
spectrum as (5.10.4).  (See Shakura and Sunyaev 1972.)  In the case of a homo\-geneous
atmosphere, eqs.\ (5.10.2) and (5.10.3) show that the ``spectrum
modification factor'' has the form
%
\begin{equation}
\tag{5.10.5}
[\tau_{es}\tau_{ff\nu*} = 1]^{-1/2}
\propto \alpha_\nu^{-3/2}(1 - e^{-x})^{1/2}
\end{equation}
%
so the spectrum has the form
%
\[
F_\nu \propto \frac{x^{3/2} e^{-x/2}}{(e^x - 1)^{1/2}}
\]
%
The total flux in this case works out to be
%
\begin{equation}
\tag{5.10.6}
F = (1.54 \times 10^{-4}~\text{erg/cm}^2~\text{sec})
\left(\frac{\rho_0 m_p}{(\text{cm}^{-3})}\right)
\left(\frac{T_s}{\text{K}}\right)^{9/4}
\end{equation}

If the atmosphere is not homogeneous, the spectrum and flux have somewhat
different forms than these.  (See Shakura and Sunyaev 1972.)

Assuming a homogeneous atmosphere with density roughly equal to that in
the central regions of the disk, and combining eq.\ (5.10.7) for the flux with
expressions (5.9.8) and (5.9.9) for the structures of the middle and inner regions,
we obtain surface temperatures of
%
\begin{align}
T_s &= (5 \times 10^7~\text{K})(\alpha^{28/80}M_*^{-1/5}\dot{M}_{0*}^{45})_*^{-87/90}\mathscr{g}^{21}\mathscr{g}^{-26/45}\,\mathscr{g}^{-2/9}
\notag\\
&\quad \text{for middle region}
\tag{5.10.8}\\[6pt]
T_s &= (6 \times 10^8~\text{K})(\alpha^{28/9}M_*^{-10/9}\dot{M}_{0*}^{8/9})_*^{-17/2}\mathscr{g}^{8/9}\,\mathscr{g}^{-16/9}\mathscr{g}^{-2/9}
\notag\\
&\quad \text{for inner region.}
\notag
\end{align}

In the inner region these surface temperatures are roughly an order of magnitude
higher than would occur if the disk could radiate as a blackbody [cf.\ eq.\
(5.2.21)].

The above estimates assume, of course, that the disk is optically thick in the
sense that $\tau_* > 1$ at the central plane $z = 0$.  However, in the innermost regions
of the disk this may not be the case.  In fact, our model [eq.\ (5.9.10)] predicts
%
\[
\tau_*(z = 0) \simeq (1 \times 10^{-3})\alpha^{-17/16}M_*^{31/16}\dot{M}_{0*}^{93/32}\,\mathscr{g}^{-25/8}\,\mathscr{g}^{41/8}
\]
%
\[
\times\,\mathscr{g}^{1/2}\,\mathscr{g}^{25/16}\,\mathscr{g}^{-2}\,.
\]

Thus, for galactic X-ray sources with $\alpha \sim 1$, $M_* \sim M_{0*} \sim 1$, $a = M$ (``maximal
Kerr''), the disk will be optically thin over a narrow region between $r \sim 2M$ and
$r \sim 10M$, outside it
will also be thick.

In such an optically thin region the disk must still radiate the huge flux
(5.6.14b) required by energy conservation.  It can do so only at temperatures far
above the blackbody value for the given flux.  All the photons emitted will
escape, so the spectrum will have the free-free form
%
\begin{equation}
\tag{5.10.9}
F_\nu \propto \alpha_{ff}^{\nu}\, e^{-x}
\end{equation}
%
(see \S\ref{sec:2.1} and \ref{sec:2.6}).  At least this would be the case if Comptonization were
negligible.  However, at such high temperatures ($T \sim 10^9$\,K), Comptonization
must be taken into account.  The free-free spectrum has many more low-energy
photons than high; and when scattered, each low-energy photon gets boosted
in energy by a fractional amount
%
\begin{equation}
\tag{5.10.10}
\Delta h\nu / h\nu \simeq 4(T/m_e) \simeq (T/2 \times 10^8~\text{K})\,.
\end{equation}

As a result the low-energy end of the spectrum gets depleted and the high-energy
end gets augmented.  [See Shakura and Sunyaev (1972) for quantitative estimates
of this effect.]

The total spectrum as observed from Earth is the integral of $F_\nu$ over the
entire disk---with corrections in the innermost regions for the capture of some
photons by the black hole.  The foundations for calculating the capture
of the radiation by the black hole are given in Bardeen's lectures in this volume.

Figure~5.10.1 shows the total spectrum as calculated by Shakura and Sunyaev
(1972) for two typical cases, without taking account of relativistic effects or
capture by the black hole.  In both cases the disk is optically thick everywhere.
In the X-ray region (1 to 100\,keV) these spectra are similar to those observed
for typical galactic X-ray sources; see lectures of Gursky in this volume.

\begin{figure}[htbp]
\centering
% \includegraphics[width=\columnwidth]{fig_5_10_1}
\fbox{\parbox{0.9\columnwidth}{\centering [Figure 5.10.1 placeholder]\\
The total power per unit frequency $L_\nu$ emitted by two model disks around
black holes, as calculated by Shakura and Sunyaev (1972) without taking account of
capture of radiation by the hole.  In both models the hole was
assumed to be nonrotating, so the inner edge of the disk was at $r_* = r/r_g = 3$.  Model (a)
corresponds to\\
$a \sim 10^{-3}$, \quad $M = M_\odot$, \quad $\dot{M}_0 = 10^{-9}\,M_\odot/\text{yr}$, \quad
$L = L_{\text{crit}} = 10^{38}~\text{erg/sec}$;\\
$\alpha \sim 10^{-2}$ to $1$ (spectrum insensitive to $\alpha$).\\
Model (b) corresponds to\\
$\alpha \sim 10^{-3}$, \quad $M = M_\odot$, \quad $\dot{M}_0 = 10^{-9}\,M_\odot/\text{yr}$,\\
$L = 10^{38}~\text{erg/sec}$;\\
$M = M_\odot$, \quad $\dot{M}_0 = 10^{-4}\,M_\odot/\text{yr}$, \quad $L = 10^{36}~\text{erg/sec}$.\\
The portion of the spectrum marked ``blackbody'' is generated primarily by the outer,
cool region of the disk where electron scattering is unimportant.  The portion marked
``modified'' is generated by the middle and inner regions where electron scattering is the
dominant source of opacity.  The temperature of the innermost region.}}
\caption{Total spectrum emitted by model accretion disks around black holes, as
calculated by Shakura and Sunyaev~\cite{ShakuraSunyaev1973}.  Curve (a) and curve (b) correspond
to two different parameter choices (see text).  The blackbody portion corresponds to the outer
region; the ``modified'' portion corresponds to the middle and inner regions dominated by
electron scattering.}
\label{fig:5.10.1}
\end{figure}

%% ================================================================
%% 5.11  Heating of Outer Region by X-Rays from the Inner Region
%% ================================================================

\subsection{Heating of the Outer Region by X-Rays from the Inner Region}
\label{sec:5.11}

Thus far we have ignored any radiative energy transfer from one portion of the
disk to another.  As long as the disk is optically thick, no significant transfer can
occur through its interior.  However, radiation can skim along the surface of the
disk from the hot, inner region toward the cooler outer region:  Because the disk
is thick with increasing radius
%
\[
h = \text{constant in inner region}
\]
\[
h \propto r^{21/20} \quad \text{in middle region}
\]
\[
h \propto r^{9/8} \quad \text{in outer region,}
\]
%
such radiation will get captured in the middle and outer regions and will heat
their surface layers.  As a result, the outer regions will thicken even more than is
predicted by the above model.  Rough estimates by Shakura and Sunyaev (1972)
suggest that for typical binary accretion models ${\sim}\,1$ per cent of the total X-ray
luminosity can be captured and reradiated.  About 10 per cent of the reradiated
flux comes off in spectral lines in the optical region, and about 90 per cent comes
off in ultraviolet free-bound and free-free emission.
